% ======================================================================
%  sec29.tex —— 第 29 节 （拟）射影态射
%  底本：原书 pp.157–160；原书标题：29. (Quasi-)projective morphisms
% ======================================================================

\section{（拟）射影态射}\label{sec:29}

\subsection{定义}\label{sec:29.1}

\begin{definition}\label{def:29.1.1}
设 $f\colon X\to Y$ 是概形间的态射。若存在有限型的拟凝聚 $\cO_Y$-模 $\cE$，使 $X$ 与 $\PP_Y(\cE)$ 的某个闭子概形 $Y$-同构，则称 $f$ 是\emph{射影的}。
\end{definition}

\begin{yizhu}
定义要求 $\cE$ 有限型拟凝聚，故 $\PP_Y(\cE)$ 在 $Y$ 上射影：在 $Y$ 的仿射开集上，$\cE$ 由有限多个截面生成，于是 $\PP_Y(\cE)$ 是该仿射开集上某个 $\PP^n$ 的闭子概形。定理 \ref{thm:29.1.5}(1) 中的紧合性正来自这一结构。
[来源：Stacks Project, Tag 01W8（Definition 29.44.1，projective morphism：$X$ 同构于某个射影丛 $\mathbf P(\cE)$ 的闭子概形，$\cE$ 拟凝聚且有限型），\url{https://stacks.math.columbia.edu/tag/01W8}（2026-11-17 访问）。]
\end{yizhu}

\begin{example}\label{ex:29.1.2}
闭浸入是射影的。
\end{example}

\begin{definition}\label{def:29.1.3}
设 $f\colon X\to Y$ 是概形间的态射。若 $f$ 有限型且存在 $f$-丰沛线丛，则称 $f$ 是\emph{拟射影的}。
\end{definition}

\begin{yizhu}
由定义，拟射影态射是有限型且带 $f$-丰沛线丛；而推论 \ref{cor:28.4.2} 说明 $f$-丰沛线丛的存在迫使 $f$ 分离。故拟射影像射自动分离，这一点在原书中由分开放置的两条结果合成，读时容易漏掉。
[依据：原书 p.147 推论 28.4.2（存在 $f$-丰沛线丛则 $f$ 分离）]
\end{yizhu}

\begin{example}\label{ex:29.1.4}
拟紧开浸入是拟射影的。更一般地，有限型的拟仿射态射是拟射影的。
\end{example}

\begin{theorem}\label{thm:29.1.5}
有：
\begin{enumerate}[label=(\arabic*),leftmargin=2.4em,itemsep=0.3ex]
  \item 射影态射 $f\colon X\to Y$ 紧合且拟射影。
  \item 若 $Y$ 拟紧且拟分离，则紧合且拟射影的态射 $f\colon X\to Y$ 是射影的。
\end{enumerate}
\end{theorem}

\begin{proof}
先证 (1)。设 $i\colon X\to\PP_Y(\cE)$ 是在 $Y$ 上定义的闭浸入，且 $\cE$ 有限型。按定义，$i^*\cE$ 是 $f$-极丰沛线丛。由推论 \ref{cor:27.3.10}，$p\colon\PP_Y(\cE)\to Y$ 紧合。故 $f=p\circ i$ 紧合。

现证 (2)。由命题 \ref{prop:28.6.2}，存在 $f$-极丰沛线丛 $\cL$。由命题 \ref{prop:28.5.6}，存在在 $Y$ 上定义的局部闭浸入 $i\colon X\to\PP_Y(\cE)$，其中 $\cE$ 有限型。因 $p\colon\PP_Y(\cE)\to Y$ 分离且 $f=p\circ i$ 紧合，态射 $i$ 紧合。故 $i$ 是闭浸入，从而 $f$ 射影，如所欲证。
\end{proof}

\begin{corollary}\label{cor:29.1.6}
设 $Y$ 是拟紧概形且存在丰沛线丛。
\begin{enumerate}[label=(\arabic*),leftmargin=2.4em,itemsep=0.3ex]
  \item 态射 $f\colon X\to Y$ 射影，当且仅当 $X$ 对某个 $n\ge 0$ 与 $\PP^n_Y$ 的闭子概形 $Y$-同构。
  \item 态射 $f\colon X\to Y$ 拟射影，当且仅当 $X$ 对某个 $n\ge 0$ 与 $\PP^n_Y$ 的拟紧局部闭子概形 $Y$-同构。
\end{enumerate}
\end{corollary}

\begin{proof}
由推论 \ref{cor:28.5.7} 得到。
\end{proof}

\begin{warning}\label{warn:29.1.7}
（拟）射影性对底概形不是局部的。事实上，Hironaka 构造了两个光滑复 $3$ 维概形之间的紧合态射 $f\colon X\to Y$，使 $f$ 在 $Y$ 的某个开覆盖上的基变换是射影的，但 $f$ 本身不是射影的。见 [Sta25, Tag 08J0]。
\end{warning}

\begin{yizhu}
Stacks 给出的具体构造是：取 $Y=\PP^3_{\C}$，$C,D$ 是相交于两点 $P,Q$ 的光滑圆锥曲线。在 $Y\setminus Q$ 上先沿 $C$ 爆破、再沿 $D$ 的严格变换爆破；在 $Y\setminus P$ 上先沿 $D$ 爆破、再沿 $C$ 的严格变换爆破。两片在 $Y\setminus\{P,Q\}$ 上典范同构，可粘成紧合而非射影的 $f$；紧合而非射影者必也非拟射影。
[来源：Stacks Project, Tag 08J0（Section 110.67，Being projective is not local on the base；Lemma 110.67.1），\url{https://stacks.math.columbia.edu/tag/08J0}（2026-11-17 访问）。]
\end{yizhu}

\subsection{Chow 引理}\label{sec:29.2}

\begin{theorem}\label{thm:29.2.1}
设 $f\colon X\to Y$ 是分离的有限型态射。假设
\begin{enumerate}[label=(\alph*),leftmargin=2.6em,itemsep=0.25ex]
  \item 概形 $Y$ 拟紧且拟分离；
  \item 概形 $X$ 只有有限多个不可约分支。
\end{enumerate}
则存在图
\begin{equation}
\begin{tikzcd}[column sep=4.4em, row sep=3em]
X & X' \arrow[l, "\pi"'] \arrow[r, "i"] & \PP^n_Y \arrow[dl] \\
& Y \arrow[ul, "f"] &
\end{tikzcd}
\label{eq:29.1}
\end{equation}
使：
\begin{enumerate}[label=(\roman*),leftmargin=2.6em,itemsep=0.25ex]
  \item 态射 $i$ 是局部闭浸入。
  \item 态射 $\pi$ 紧合且满。
  \item 存在稠密开子集 $U\subseteq X$，使 $\pi^{-1}(U)$ 在 $X'$ 中稠密，且 $\pi^{-1}(U)\to U$ 是同构。
\end{enumerate}
\end{theorem}

\begin{remark}\label{rem:29.2.2}
若 $Y$ 诺特，则定理 \ref{thm:29.2.1} 的假设 (a) 与 (b) 自动满足。

去掉假设 (b) 后，仍存在只满足 (i) 与 (ii) 的图 \eqref{eq:29.1}。这可借助诺特逼近证明。见 [Sta25, Tag 01ZZ]。
\end{remark}

\begin{yizhu}
Stacks 在所引的变体里给出：$S$ 拟紧拟分离、$f\colon X\to S$ 分离有限型时，存在 $n\ge 0$ 与图 $X\xleftarrow{\pi}X'\xrightarrow{i}\PP^n_S$，其中 $i$ 是浸入、$\pi$ 紧合且满。证明写成 $S=\lim S_i$ 的诺特逼近，把问题化到诺特底概形的情形再用定理 \ref{thm:29.2.1}。可见去掉 (b) 时丢掉的正是 (iii) 这条「在稠密开集上为同构」。
[来源：Stacks Project, Tag 01ZZ（Section 32.12，Variants of Chow's Lemma；Lemma 32.12.1），\url{https://stacks.math.columbia.edu/tag/01ZZ}（2026-11-17 访问）。]
\end{yizhu}

\begin{remark}\label{rem:29.2.3}
在定理 \ref{thm:29.2.1} 中，态射 $\pi$ 实为射影的。事实上，$X'$ 与 $\PP^n_X$ 的某个闭子概形 $X$-同构。

为看出这点，注意局部闭浸入 $i$ 分解为
\[
X'\xrightarrow{g}\PP^n_Y\times_Y X\to\PP^n_Y,
\]
故 $g$ 是局部闭浸入。同样，紧合态射 $\pi$ 分解为
\[
X'\xrightarrow{g}\PP^n_Y\times_Y X\to X,
\]
故 $g$ 紧合。合起来看，$g$ 是闭浸入，如所欲证。
\end{remark}

\begin{exercise}\label{exc:29.2.4}
在定理 \ref{thm:29.2.1} 中，若 $f$ 紧合，则 $i$ 是闭浸入。
\end{exercise}

\begin{remark}\label{rem:29.2.5}
特别地，定理 \ref{thm:29.2.1} 说明：任意在 $Y$ 上紧合的概形 $X$，都双有理等价于某个在 $Y$ 上射影的概形 $X'$。这正是复几何中 Chow 引理的原初陈述。
\end{remark}

\begin{proof}[定理 \ref{thm:29.2.1} 的证明]
先注意可以把 (iii) 换成
\begin{enumerate}[label=(\roman*),leftmargin=2.6em,itemsep=0.25ex,start=3]
  \item[\textup{(iii')}] 存在稠密开子集 $U\subseteq X$，使 $\pi^{-1}(U)\to U$ 是同构。
\end{enumerate}
即：设已有满足 (i)、(ii) 与 (iii') 的图 \eqref{eq:29.1}，则把 $X'$ 换成 $\pi^{-1}(U)$ 在 $X'$ 中的概形论闭包，所得图仍满足 (i)、(ii) 与 (iii)。

注意 (a) 蕴含 $X$ 拟紧。

\begin{exercise}\label{exc:29.2.6}
设 $X$ 是拟紧概形且只有有限多个不可约分支。证明：存在由仿射开子概形组成的有限覆盖 $X=\bigcup_{i=1}^nU_i$，使 $U:=\bigcap_{i=1}^nU_i$ 在 $X$ 中稠密。
\end{exercise}

取如上述习题中的 $U_i$ 与 $U$。因每个 $U_i$ 仿射且在 $Y$ 上有限型，由引理 \ref{lem:28.6.3}（施于 $U_i\to Y$ 与 $\cO_{U_i}$），存在局部闭浸入 $\iota_i\colon U_i\to\PP^{m_i}_Y$（某个 $m_i\ge 0$）。注意由假设 (a)，$\iota_i$ 拟紧。又
\[
\iota\colon U\to U\times_Y\cdots\times_Y U\xrightarrow{(\iota_1,\cdots,\iota_n)}\PP^{m_1}_Y\times_Y\cdots\times_Y\PP^{m_n}_Y=:P
\]
也是拟紧局部闭浸入。令 $Z_i:=\overline{U_i}$、$Z:=\overline U$ 为这些浸入的概形论闭包。回忆概形论闭包的构造具有函子性（注记 \ref{rem:18.5.3}）。故对每个 $i$，有交换图
\[
\begin{tikzcd}[column sep=3.6em, row sep=2.6em]
& U \arrow[dl] \arrow[d] \arrow[dr] & \\
V_i \arrow[r] \arrow[d, "q_i"'] & Z \arrow[r, "\bar\iota"] \arrow[d, "p_i"] & P \arrow[d, "pr_i"] \\
U_i \arrow[r] & Z_i \arrow[r, "\bar\iota_i"'] & \PP^{m_i}_Y,
\end{tikzcd}
\]
其中
\[
V_i := U_i\times_{Z_i}Z\simeq p_i^{-1}(U_i)
\]
是 $Z$ 的含 $U$ 的开子概形。注意 $\iota_i$、$\iota$ 与 $pr_i$ 都紧合。故 $p_i$ 紧合，其基变换 $q_i$ 也紧合。

现定义
\[
X' := \bigcup_{i=1}^nV_i \subseteq Z
\]
为这些开子概形之并。注意
\[
X'\to Z\xrightarrow{\bar\iota}P\to\PP^m_Y
\]
确实是局部闭浸入，其中最后一个态射是 Segre 嵌入
\[
\PP^{m_1}_Y\times_Y\cdots\times_Y\PP^{m_n}_Y\to\PP^m_Y,\qquad m:=\prod_{i=1}^n(m_i+1)-1.
\]

为完成证明，还需构造 $Y$-态射 $\pi\colon X'\to X$，使：
\begin{enumerate}[label=(\arabic*),leftmargin=2.4em,itemsep=0.25ex]
  \item 态射 $\pi$ 紧合；
  \item 态射 $\pi$ 满；
  \item 态射 $\pi$ 诱导同构 $\pi^{-1}(U)\to U$。
\end{enumerate}

为此先证：存在唯一的 $Y$-态射 $\pi\colon X'\to X$，使 $\pi|_{V_i}$ 由复合
\[
\pi_i\colon V_i\xrightarrow{q_i}U_i\xrightarrow{\subseteq}X
\]
给出。只需验证 $\pi_i|_{V_i\cap V_j}=\pi_j|_{V_i\cap V_j}$。因 $Z$ 是 $U\to P$ 的概形论闭包，它也是 $U\to Z$ 的概形论闭包。由此 $V_i\cap V_j$ 是 $U\to V_i\cap V_j$ 的概形论闭包。由构造，$\pi_i|_{V_i\cap V_j}$ 与 $\pi_j|_{V_i\cap V_j}$ 在 $U$ 上的限制相等，因为二者都由典范开浸入 $U\to X$ 给出。故 $\pi_i|_{V_i\cap V_j}$ 与 $\pi_j|_{V_i\cap V_j}$ 相等，因为 $X$ 在 $Y$ 上分离。

现在证明如上构造的态射 $\pi$ 满足 $\pi^{-1}(U_i)=V_i$。显然 $V_i$ 是 $\pi^{-1}(U_i)$ 的开子概形。注意复合
\[
\pi^{-1}(U_i)\to U_i\to Y
\]
分离，因为 $\pi^{-1}(U_i)$ 是 $P$ 的局部闭子概形，而 $P$ 在 $Y$ 上分离。故 $\pi^{-1}(U_i)\to U_i$ 分离。又复合
\[
q\colon V_i\to\pi^{-1}(U_i)\to U_i
\]
紧合，特别地闭，故 $V_i\to\pi^{-1}(U_i)$ 也闭。于是 $V_i=\pi^{-1}(U_i)$。

态射 $\pi$ 紧合，因为它到 $U_i\subseteq X$ 的基变换由紧合态射 $q_i\colon V_i\to U_i$ 给出。此即 (1)。

态射 $\pi$ 满，因为它紧合、特别地闭，且其像含 $U$。此即 (2)。

剩下证 (3)。注意恒等态射 $\id_U$ 分解为
\[
\id_U\colon U\to\pi^{-1}(U)\to U,
\]
其中第一个态射是（$X'$ 的开子概形之间的）开浸入。因复合与第二个态射都紧合，第一个态射也紧合。这说明 $U=\pi^{-1}(U)$，如所欲证。
\end{proof}

\begin{yizhu}
证明中「取有限仿射覆盖使交稠密」这一步（习题 \ref{exc:29.2.6}）是为造出公共的 $U$：$\iota$ 定义在 $U$ 上，$Z$ 取 $U$ 的闭包，$X'$ 再由各 $V_i=U_i\times_{Z_i}Z$ 拼成。若 $X$ 不可约，任取非空仿射开集即可；一般情形必须先去掉多余分支，否则 $\bigcap U_i$ 可能空，$Z$ 也就无从谈起。
[译者按（上下文推断，无外部依据）：此处说明习题 \ref{exc:29.2.6} 在整条证明中的作用]
\end{yizhu}
